\documentclass[../main.tex]{subfiles}

\begin{document}

% ============================================================
\chapter{State of the Art}
\label{chap:related}
% ============================================================

The machine learning pipeline is commonly conceptualized as a sequence of stages: data collection, data preparation, feature engineering, model training, evaluation, and deployment. While the research community has historically devoted the greatest attention to the model training stage, designing novel architectures, loss functions, and optimization algorithms (model centric AI), there is a growing body of evidence suggests that data quality issues originating \emph{upstream} of model training often have a larger impact on final model performance than the choice of model itself.

The emergence of the \emph{Data-Centric AI} movement formalizes this insight: rather than
treating the dataset as a fixed input and iterating on the model, Data-Centric
AI advocates for systematically iterating on the data: cleaning, augmenting,
curating, and validating it; while keeping the model architecture constant.
Therefore this paradigm has \emph{data quality} at the center of the ML engineering workflow. This motivates the increasing development of tools that can measure, diagnose, and, for experimental purposes, deliberately manipulate data quality dimensions. 
Those tools (such as PuckTrick) can be applied to a wide range of domains 
(such as cybersecurity, in our case). 


% ============================================================
\newpage
\section{State of the Art for Pucktrick Library}
% ============================================================

This chapter provides a detailed description of both the main form of data quality errors this thesis addresses and the PuckTrick library as it existed \emph{prior} to the modifications introduced in this work.
This corresponds to the Pandas-only version documented in the ICAIF 2025 paper (~\supercite{icaif2025pucktrick}): establishing the baseline against which the contributions of this thesis are measured and explainig why it bacame necessary for PuckTrick to be upgraded.

% ============================================================
\subsection{Limitations of the Pre-Extension PuckTrick}
\label{sec:pucktrick-limitations}
% ============================================================

PuckTrick distinguishes itself from many other available tools (such as DataAssert, CleanML, etc..) by providing a unified, modular, and extensible API for the \emph{deliberate injection} of five
orthogonal categories of data corruption, each with fine-grained control over
the corruption percentage, the target column, and the variant (New or Extended).
Its focus on controlled degradation, rather than detection or cleaning. This makes PuckTrick fill a very specific niche in the Data-Centric AI toolboxes scene. 

The library's design is guided by following principles:
\begin{enumerate}
    \item \textbf{Controlled contamination}: every corruption operation is
          parametrised by a target percentage, a target column, and (where
          appropriate) a reference to the original clean dataset, ensuring
          that the exact degree of corruption is known and reproducible.
    \item \textbf{Method-specific specialisation}: each corruption category
          (missing, noise, outliers, labels, duplicates) provides sub-methods
          tailored to the data type of the target column (continuous, discrete,
          binary, categorical string, categorical integer), because the
          semantics of corruption differ fundamentally across data types.
    \item \textbf{New vs.\ Extended paradigm}: every corruption method exists
          in two variants:
          \begin{itemize}
              \item \emph{New} methods inject corruption into a \emph{clean}
                    column (one that has not been previously corrupted), reaching
                    the specified percentage from zero.
              \item \emph{Extended} methods take a column that has already been
                    partially corrupted and add further corruption until the
                    cumulative percentage reaches the specified target.  This
                    two-phase design supports incremental experiments in which
                    corruption is progressively increased.
          \end{itemize}

    \newpage
    \item \textbf{Wide range of usable methods}: PuckTrick offers a total of 5 method to insert noise inside datasets. Table~\ref{tab:dimensions} provides a summary:

    \begin{table}[H]
    \centering
    \caption{Data quality error dimensions implemented by
      \textsc{PuckTrick}.}
    \label{tab:dimensions}
    \small
    \begin{tabularx}{\textwidth}{llX}
      \toprule
      \textbf{Dimension} & \textbf{Module} & \textbf{Description} \\
      \midrule
      Duplicated records & \texttt{duplicated.py} &
        Introduces duplicate rows. \\ \\
      Label noise & \texttt{labels.py} &
        Flips binary labels ($0 \leftrightarrow 1$) or replaces
        categorical target values with different existing categories. \\ \\
      Missing values & \texttt{missing.py} &
        Sets selected cell values to \texttt{NaN} (Pandas)\\ \\
      Random noise & \texttt{noisy.py} &
        Replaces column values with random alternatives drawn from the column's observed range $[\min, \max]
        \newline
        $(strings, numbers, dates, or categories). \\ \\
      Outliers & \texttt{outliers.py} &
        Injects extreme values beyond the normal range: for numeric
        columns, values beyond $\mu \pm 4\sigma$; for categorical
        integers, the next integer beyond the maximum; for strings,
        a sentinel value; for dates, extreme timestamps. \\
      \bottomrule
    \end{tabularx}
    \end{table}

\end{enumerate}

The data flow is straightforward: the user passes a Pandas DataFrame and a
\emph{strategy dictionary} specifying: which corruption methods to apply, to
which columns, and at which percentages.  PuckTrick returns a new DataFrame
containing the corrupted data, leaving the original DataFrame unmodified.


While PuckTrick version 0.5.0 provides a comprehensive and well-designed
corruption toolkit, it is subject to several limitations that motivated the
extension work described in this thesis:

\begin{enumerate}
    \item \textbf{Single-machine scalability.}  All corruption operations are
          implemented on Pandas DataFrames, which reside entirely in the
          memory of a single machine.  For datasets exceeding the available
          RAM (such as the 16-million-row CIC-IDS2018 dataset used in this
          work), Pandas becomes a bottleneck, either requiring expensive
          vertical scaling or forcing the user to work on subsampled data.

    \item \textbf{No distributed computing support.}  There is no facility
          for distributing corruption operations across multiple nodes.  In
          environments where the data already resides in a distributed file
          system (e.g., HDFS, S3) or is managed by a distributed processing
          engine (e.g., Spark), the user must first collect the data to a
          single node, apply PuckTrick, and then redistribute the result: an
          inefficient and error-prone workflow.

    \item \textbf{Index-dependent implementations.}  Many corruption
          algorithms rely on Pandas's integer-positional indexing
          (\texttt{.iloc}) to select target rows, a mechanism that has no
          direct analogue in Spark's distributed DataFrame API.  This
          architectural coupling to Pandas's row-indexed data model
          makes migration to a distributed backend non-trivial. This increases necessity for PuckTrick library to have a native spark implementation, that would otherwise be impossible to use well optimized Spark operations. 

    \item \textbf{In-place mutation patterns.}  Some implementations modify
          the DataFrame in place (e.g., assigning \texttt{NaN} to specific
          cells via \texttt{.loc}), which conflicts with Spark's immutable
          DataFrame semantics where every transformation produces a new
          DataFrame.
\end{enumerate}

These limitations collectively prevented the use of PuckTrick on the full
16-million-instance CIC-IDS2018 dataset within a distributed Spark cluster: 
a capability that is essential for the large-scale noise-injection experiments that we'll explain later on in this thesis (Chapter \ref{chap:Empiric-analysis-models} and \ref{chap:Empiric-analysis-models}).  

Therefore, the choice of Apache Spark as the target distributed backend for PuckTrick was
motivated by its maturity, its rich DataFrame API (which most closely mirrors
the operations required by PuckTrick), its native integration with distributed
storage systems (HDFS, S3, Parquet), and its widespread adoption in both
industry and academic data engineering pipelines.






% ============================================================
\newpage
\section{Deep Learning for NIDS}
\label{sec:nids-dl}
% ============================================================

A second object of this thesis is studying how NIDSs systems can deteriorate or either be improved by using systematical noise injection. 


% ============================================================
\subsection{The Impact of Data Quality on Model Performance}
\label{sec:dq-impact}
% ============================================================

Before this thesis's work, a substantial body of empirical research has investigated the sensitivity of machine learning models to the aforementioned data quality issues. 
If we take a look at the general results obtained while using same type of noise as we will use when applying PuckTricks methods we get the following results: 

\paragraph{Missing values.}
The effect of missing data depends on the missingness mechanism, the proportion
of missing values, and the imputation strategy employed.
Garc\'{\i}a-Laencina \emph{et al.}~\supercite{garciaLaencina2010pattern} demonstrate
that even moderate rates of missingness (5 to 10\%) can significantly degrade
classifier accuracy, especially for algorithms that cannot handle missing inputs
natively (e.g., standard neural networks).  Tree-based ensembles (e.g., XGBoost,
Random Forest) and TabNet are more robust due to their ability to route
instances around missing splits or to learn attention masks that down-weight
incomplete features.

\paragraph{Feature noise.}
Zhu and Wu~\supercite{zhu2004class} show that feature noise generally causes a
monotonic decrease in model performance as the noise rate increases, with the
magnitude of the effect being highly dependent on the feature's importance to
the classification task.  Noise applied to irrelevant or redundant features has
negligible impact, while noise applied to the most discriminative features can
cause catastrophic degradation: a finding that underpins the \emph{fragility-versus-relevance} analysis performed in this thesis.

\paragraph{Label noise.}
Label noise is widely regarded as more damaging than feature
noise~\supercite{frenay2014classification, nettleton2010study}.  Even a 5\% rate of
random label flipping can reduce classifier accuracy by several percentage
points, while systematic (class-conditional) label noise can introduce
consistent bias that is especially difficult to detect.  Deep neural networks
are particularly susceptible because of their capacity to memorise noisy
labels~\supercite{arpit2017closer, zhang2021understanding}.

\paragraph{Duplicates.}
Elmagarmid \emph{et al.}~\supercite{elmagarmid2007duplicate} survey the problem of
duplicate detection and its impact on data analytics.  In the ML context,
duplicates that span training and test partitions violate the independence
assumption and lead to over-optimistic performance estimates.  Even when
duplicates are confined to the training set, they bias the model towards
over-represented instances, effectively functioning as a form of uncontrolled
oversampling that can exacerbate class imbalance effects.

\paragraph{Outliers.}
Outliers exert a disproportionate influence on models that minimize
squared-error losses or that are sensitive to feature scaling (e.g., linear
models, $k$-NN, SVMs with RBF kernels).  Deep learning models with batch
normalization layers exhibit some degree of natural robustness to outliers, but
extreme values in the input may still cause gradient instability, especially
during the initial phases of training (~\supercite{goodfellow2016deep}).


% ============================================================
\subsection{State of the Art for Deep Learning on NIDS}
\label{sec:dq-impact}
% ============================================================

Network Intrusion Detection Systems (NIDS) aim to identify malicious network activity by analysing traffic metadata and, in some cases, packet payloads. Traditional NIDS relied on signature-based matching.  Yet, ith the advent of deep learning,
data-driven approaches have achieved substantially higher detection rates,
especially for zero-day attacks that lack known signatures.

The CIC-IDS2018 dataset (~\supercite{cicids2018}), used as the experimental testbed in
this thesis, has been widely adopted in the NIDS literature.  Prior works have
applied a variety of ML and DL architectures to this dataset:

\begin{itemize}
    \item \textbf{Random Forest and Gradient Boosting} models achieve strong
          baseline performance on flow-level features, with reported F1 scores
          exceeding 0.95 on balanced subsets (~\supercite{sharafaldin2018toward}).
    \item \textbf{Convolutional Neural Networks (CNNs)} have been applied to
          reshaped flow vectors or to traffic image representations, exploiting
          spatial correlations between features ~\supercite{kim2020cnn}.
    \item \textbf{LSTM and GRU networks} capture temporal dependencies in
          sequential flow data, improving detection of multi-step attacks such
          as infiltration and slow-rate DoS (~\supercite{imrana2021bidirectional}).
    \item \textbf{Autoencoders and Variational Autoencoders} have been used
          for unsupervised anomaly detection, learning a compressed
          representation of benign traffic and flagging flows that cannot be
          accurately reconstructed (~\supercite{mirsky2018kitsune}).
\end{itemize}

However, none of the above works systematically investigate how \textbf{data
quality degradation} affects model performance.  The typical experimental
protocol assumes a clean dataset (or applies ad-hoc cleaning) and evaluates
model accuracy on a held-out test set.  This thesis \textbf{fills this gap} by
introducing controlled corruptions via PuckTrick and measuring their impact on
TabNet and a specialized CNN-LSTM in a multi-task/multi-target configuration. 
This could then produce \textbf{real word applications} on what type of noise could be most dangerous for this type of task and at the same time could also lead to discovery of how systematic noise injection for this task could increase NIDS models resistance to identifying cyber-attacks.

\end{document}